\Needspace{8\baselineskip}
\Article{197}
Haben wir eine Transformation einer Form $(a,b,c)\mathinner{\ldotp\ldotp\ldotp}\varphi$ in eine äquivalente Form $\Phi$, so können wir daraus \textit{alle} gleichartigen Transformationen von $\varphi$ in $\Phi$ ableiten, sofern wir alle Lösungen der unbestimmten Gleichung $tt-Duu=mm$ angeben
\SourceLeaf{192}
können. Dabei bezeichnet \(D\) die Determinante der Formen \(\Phi,\varphi\) und \(m\) den größten gemeinsamen Teiler von \(a,2b,c\) (Artikel~162). Dieses Problem haben wir oben bereits für einen negativen Wert von \(D\) gelöst; nun gehen wir es für einen positiven Wert an. Offenbar erfüllt jeder Wert von \(t\), der der Gleichung genügt, sie auch nach Umkehrung seines Vorzeichens; dasselbe gilt für jeden Wert von \(u\). Daher genügt es, alle \textit{positiven} Werte von \(t,u\) anzugeben, und jede Lösung mit positiven Werten vertritt vier Lösungen. Wir werden dazu \textit{zuerst} zeigen, wie man die \textit{kleinsten} Werte von \(t,u\) (außer den unmittelbar ersichtlichen \(t=m,\ u=0\)) findet, und \textit{dann}, wie man daraus alle übrigen ableitet.

\Needspace{8\baselineskip}
\Article{198}
\textsc{Aufgabe.} \textit{Die kleinsten Zahlen \(t,u\) finden, die der unbestimmten Gleichung \(tt-Duu=mm\) genügen, wenn eine Form \((M,N,P)\) mit der Determinante \(D\) gegeben ist und die Zahlen \(M,2N,P\) den größten gemeinsamen Teiler \(n\) haben.}

\textit{Lösung.} Wir wählen eine beliebige reduzierte Form \((a,b,-a')\ldots f\) mit der Determinante \(D\), bei der die Zahlen \(a,2b,a'\) den größten gemeinsamen Teiler \(m\) haben. Dass es eine solche gibt, ist schon daraus klar, dass wir eine zu \((M,N,P)\) äquivalente reduzierte Form finden können, die nach Artikel~161 diese Eigenschaft besitzt. Für unseren jetzigen Zweck können wir aber jede reduzierte Form verwenden, die diese Bedingung erfüllt. Wir bestimmen die Periode von \(f\), die aus \(n\) Formen bestehen möge. Mit allen Bezeichnungen aus Artikel~188 ist \(f^n\) gleich \((+a^n,b^n,-a^{n+1})\), weil \(n\) gerade ist; in diese Form geht \(f\) durch die eigentliche Substitution \(\alpha^n,\beta^n,\gamma^n,\delta^n\) über. Da aber \(f\) und \(f^n\) identisch sind, geht \(f\) auch in \(f^n\) über, und zwar durch die eigentliche Substitution \(1,0,0,1\). Aus diesen beiden gleichartigen Transformationen von \(f\) in \(f^n\) können wir nach Artikel~162 eine Lösung der Gleichung \(tt-Duu=mm\) \textit{in ganzen Zahlen} ableiten, nämlich \(t=\frac12(\alpha^n+\delta^n)m\) (Gleichung~18, Artikel~162), \(u=\frac{\gamma^n m}{a}\) (Gleichung~19)\HistoricalNote{*)}{\begin{tabular}[t]{@{}l@{}}\(\begin{array}{@{}rll@{}}
\text{In Artikel 162 hatten wir}&\alpha,\beta,\gamma,\delta;&\alpha',\beta',\gamma',\delta';\\
\text{hier haben wir}&1,0,0,1;&\alpha^n,\beta^n,\gamma^n,\delta^n;\end{array}\)
\\[0.6ex]\(\begin{array}{@{}lll@{}}A,B,C;&a,b,c;&e,\\a,b,-a';&a,b,-a';&1.\end{array}\)\end{tabular}}. Wir bezeichnen diese Werte, positiv genommen, falls sie es noch nicht sind, mit \(T,U\). Dann sind \(T,U\) die kleinsten Werte von \(t,u\), außer \(t=m,\ u=0\) (von diesen müssen sie verschieden sein, da \(\gamma^n\) offenbar nicht \(=0\) sein kann).

Nehmen wir nämlich an, es gebe noch kleinere Werte von \(t,u\), etwa \(\mathfrak{t},\mathfrak{u}\), die positiv sind, wobei \(\mathfrak{u}\) nicht \(=0\) ist. Dann wird nach Artikel~162 die Form \(f\) durch die eigentliche Substitution \(\frac1m(\mathfrak{t}-b\mathfrak{u}),\ \frac1m a'\mathfrak{u},\ \frac1m a\mathfrak{u},\ \frac1m(\mathfrak{t}+b\mathfrak{u})\) in eine mit ihr
\SourceLeaf{193}
identische Form transformiert. Nach Artikel~193, II muss \textit{entweder} \(\frac1m(\mathfrak{t}-b\mathfrak{u})\) \textit{oder} \(-\frac1m(\mathfrak{t}-b\mathfrak{u})\) einer der Zahlen \(\alpha'',\alpha''',\alpha''''\) usw. gleich sein, etwa \(=\alpha^\mu\) (denn wegen \(\mathfrak{tt}=D\mathfrak{uu}+mm=bb\mathfrak{uu}+aa'\mathfrak{uu}+mm\) ist \(\mathfrak{tt}>bb\mathfrak{uu}\) und daher \(\mathfrak{t}-b\mathfrak{u}\) positiv; somit hat der Bruch \(\frac{\mathfrak{t}-b\mathfrak{u}}{a\mathfrak{u}}\), der dem Bruch \(\frac{\mathfrak{A}}{\mathfrak{C}}\) in Artikel~193 entspricht, dasselbe Vorzeichen wie \(a\) oder \(a'\)). Im ersten Fall sind \(\frac1m a'\mathfrak{u},\ \frac1m a\mathfrak{u},\ \frac1m(\mathfrak{t}+b\mathfrak{u})\), im zweiten dieselben Größen mit umgekehrten Vorzeichen der Reihe nach \(=\beta^\mu,\gamma^\mu,\delta^\mu\). Da aber \(\mathfrak{u}<U\), \textit{d.h.}, \(\mathfrak{u}<\frac{\gamma^n m}{a}\) und \(>0\) ist, gilt \(\gamma^\mu<\gamma^n\) und \(>0\). Weil die Folge \(\gamma,\gamma',\gamma''\) usw. ständig wächst, muss \(\mu\) strikt zwischen \(0\) und \(n\) liegen. Die entsprechende Form \(f^\mu\) ist aber mit \(f\) identisch, \textit{ein Widerspruch}, da alle Formen \(f,f',f''\) usw. bis \(f^{n-1}\) als verschieden vorausgesetzt sind. Also sind die kleinsten Werte von \(t,u\) (außer \(m,0\)) gleich \(T,U\).

\textit{Beispiel.} Für \(D=79,\ m=1\) können wir die Form \((3,8,-5)\) verwenden, für die \(n=6\) und \(\alpha^n=-8,\ \gamma^n=-27,\ \delta^n=-152\) gilt (Artikel~188). Daraus folgt \(T=80,\ U=9\), die kleinsten Werte von \(t,u\), die der Gleichung \(tt-79uu=1\) genügen.

\Needspace{8\baselineskip}
\Article{199}
Für die praktische Rechnung lassen sich noch bequemere Formeln finden. Es gilt nämlich \(2b\gamma^n=-a(\alpha^n-\delta^n)\), was sich leicht aus Artikel~162 ergibt, indem wir Gleichung~\([19]\) mit \(2b\) und \([20]\) mit \(a\) multiplizieren und die dortigen Bezeichnungen durch die jetzigen ersetzen. Daraus folgt \(\alpha^n+\delta^n=2\delta^n-\frac{2b}{a}\gamma^n\) und damit
\[
\pm T=m\left(\delta^n-\frac{b}{a}\gamma^n\right),\quad \pm U=\frac{\gamma^n m}{a}
\]
Auf ähnliche Weise erhalten wir
\[
\pm T=m\left(\alpha^n+\frac{b}{a'}\beta^n\right),\quad \pm U=\frac{\beta^n m}{a'}
\]
Beide Formelpaare sind sehr bequem, weil \(\gamma^n=\delta^{n-1}\), \(\alpha^n=\beta^{n-1}\) ist. Verwenden wir das zweite, so brauchen wir nur die Folge \(\beta',\beta'',\beta'''\ldots\beta^n\) zu berechnen; bevorzugen wir das erste, nur die Folge \(\delta',\delta'',\delta'''\) usw. Außerdem folgt leicht aus Artikel~189, 3, dass, weil \(n\) notwendig gerade ist, \(\alpha^n\) und \(\frac{b}{a'}\beta^n\) dasselbe Vorzeichen haben, ebenso \(\delta^n\) und \(\frac{b}{a}\gamma^n\). Daher müssen wir in der ersten Formel für \(T\) die absolute Differenz und in der zweiten die absolute Summe nehmen und brauchen die Vorzeichen überhaupt nicht zu beachten. Mit den Bezeichnungen aus Artikel~189, 4 ergibt die erste Formel
\[
\begin{aligned}T&=m[k',k'',k'''\ldotp\,\ldotp\,\ldotp\,\ldotp\, k^n]-\frac{mb}{a}[k',k'',k'''\ldotp\,\ldotp\,\ldotp\,\ldotp\, k^{n-1}],\\U&=\frac{m}{a}[k',k'',k'''\ldotp\,\ldotp\,\ldotp\,\ldotp\, k^{n-1}]\end{aligned}
\]
\SourceLeaf{194}
und die zweite
\[
\begin{aligned}T&=m[k'',k'''\ldotp\,\ldotp\,\ldotp\,\ldotp\, k^{n-1}]+\frac{mb}{a'}[k'',k'''\ldotp\,\ldotp\,\ldotp\,\ldotp\, k^n],\\U&=\frac{m}{a'}[k'',k'''\ldotp\,\ldotp\,\ldotp\,\ldotp\, k^n]\end{aligned}
\]
wobei wir für den Wert von \(T\) auch \(m\left[k'',k'''\ldotp\,\ldotp\,\ldotp\,\ldotp\, k^n,\frac{b}{a'}\right]\) schreiben können.

\textit{Beispiel.} Für \(D=61,\ m=2\) können wir die Form \((2,7,-6)\) verwenden, für die wir \(n=6\) und \(k',k'',k''',k'''',k''''',k''''''\) der Reihe nach \(=2,2,7,2,2,7\) finden. Daraus folgt
\[
\begin{aligned}T&=2[2,2,7,2,2,7]-7[2,2,7,2,2]\\&=2888-1365=1523\end{aligned}
\]
nach der ersten Formel; die zweite ergibt dasselbe:
\[
T=2[2,7,2,2]+\tfrac73[2,7,2,2.7]
\]
\[
U\text{ wird }=[2,2,7,2,2]=\tfrac13[2,7,2,2,7]=195.
\]
Es gibt noch weitere Möglichkeiten, die Rechnung zu verkürzen; um uns kurz zu halten, können wir sie hier aber nicht ausführlicher behandeln.

\Needspace{8\baselineskip}
\Article{200}
Um aus den kleinsten Werten von \(t,u\) \textit{alle} zu erhalten, schreiben wir die Gleichung \(TT-DUU=mm\) in der Form
\[
\left(\frac{T}{m}+\frac{U}{m}\surd D\right)\left(\frac{T}{m}-\frac{U}{m}\surd D\right)=1
\]
woraus auch folgt
\[
\left(\frac{T}{m}+\frac{U}{m}\surd D\right)^e\left(\frac{T}{m}-\frac{U}{m}\surd D\right)^e=1\quad\mathinner{\ldotp\ldotp\ldotp\ldotp\ldotp\ldotp\ldotp}\tag*{[1]}
\]
wobei \(e\) eine beliebige Zahl bezeichnet. Wir bezeichnen nun der Kürze wegen die Werte der Größen
\begingroup\postdisplaypenalty=10000
\[
\begin{gathered}\frac{m}{2}\left(\frac{T}{m}+\frac{U}{m}\surd D\right)^e+\frac{m}{2}\left(\frac{T}{m}-\frac{U}{m}\surd D\right)^e,\\
\frac{m}{2\surd D}\left(\frac{T}{m}+\frac{U}{m}\surd D\right)^e-\frac{m}{2\surd D}\left(\frac{T}{m}-\frac{U}{m}\surd D\right)^e\text{\textsuperscript{*)}}\end{gathered}
\]
\HistoricalDisplayNote{*)}{Nur in diesen vier Ausdrücken und in Gleichung~\([1]\) bezeichnet \(e\) den \textit{Exponenten} einer Potenz; sonst bezeichnen hochgestellte Buchstaben stets einen \textit{Index}.}\endgroup
allgemein mit \(t^e,u^e\) der Reihe nach, \textit{d.h.} ihre Werte für \(e=0\) mit \(t^0,u^0\) (die gleich \(m,0\) sind); für \(e=1\) mit \(t',u'\) (die gleich \(T,U\) sind); für \(e=2\) mit \(t'',u''\); für \(t=3\) mit \(t''',u'''\) usw. --- Wir werden beweisen, dass diese Ausdrücke, wenn wir für \(e\) alle nichtnegativen ganzen Zahlen nehmen, \textit{d.h.} \(0\) und alle positiven Zahlen von \(1\) bis \(\infty\), sämtliche positiven Werte von \(t,u\) darstellen: nämlich I) dass alle Werte dieser Ausdrücke
\SourceLeaf{195}
tatsächlich Werte von \(t,u\) sind; II) dass alle diese Werte ganze Zahlen sind; III) dass es keine positiven Werte von \(t,u\) gibt, die nicht in diesen Formeln enthalten sind.

I. Setzen wir für \(t^e,u^e\) ihre Werte ein, so bestätigen wir mithilfe von Gleichung~\([1]\) leicht, dass
\[
(t^e+u^e\surd D)(t^e-u^e\surd D)=mm,\quad \textit{d.h.}\quad t^e t^e-Du^e u^e=mm
\]

II. Auf dieselbe Weise bestätigen wir leicht, dass allgemein
\[
t^{e+1}+t^{e-1}=\frac{2T}{m}t^e,\quad u^{e+1}+u^{e-1}=\frac{2T}{m}u^e
\]
Die beiden Folgen \(t^0,t',t'',t'''\) usw., \(u^0,u',u'',u'''\) usw. erfüllen also Rekursionen, beide mit den Koeffizienten \(\frac{2T}{m},-1\), nämlich
\[
t''=\frac{2T}{m}t'-t^0,\quad t'''=\frac{2T}{m}t''-t'\text{ usw.,}\quad u''=\frac{2T}{m}u'\text{ usw.}
\]
Nach Voraussetzung gibt es nun eine Form \((M,N,P)\) mit der Determinante \(D\), in der \(M,2N,P\) durch \(m\) teilbar sind. Daher gilt
\[
TT=(NN-MP)UU+mm
\]
und somit ist \(4TT\) offenbar durch \(mm\) teilbar. Also ist \(\frac{2T}{m}\) eine ganze und zwar positive Zahl. Es gilt aber \(t^0=m,\ t'=T,\ u^0=0,\ u'=U\); diese Werte sind also ganzzahlig. Damit sind auch alle \(t'',t'''\) usw., \(u'',u'''\) usw. ganzzahlig. Ferner sind wegen \(TT>mm\) offenbar alle \(t^0,t',t'',t'''\) usw. positiv und wachsen ständig unbegrenzt an, ebenso alle \(u^0,u',u'',u'''\) usw.

III. Nehmen wir an, es gebe noch andere positive Werte von \(t,u\), die nicht in den Folgen \(t^0,t',t''\) usw., \(u^0,u',u''\) usw. enthalten sind, etwa \(\mathfrak{T},\mathfrak{U}\). Da die Folge \(u^0,u'\) usw. von \(0\) aus unbegrenzt wächst, muss \(\mathfrak{U}\) offenbar zwischen zwei benachbarten Gliedern \(u^n\) und \(u^{n+1}\) liegen, sodass \(\mathfrak{U}>u^n\) und \(\mathfrak{U}<u^{n+1}\) ist. Um zu zeigen, dass diese Annahme zu einem Widerspruch führt, bemerken wir Folgendes.

\(1^\circ\) Die Gleichung \(tt-Duu=mm\) wird auch erfüllt, wenn wir setzen
\[
t=\frac1m(\mathfrak{T}t^n-D\mathfrak{U}u^n),\quad u=\frac1m(\mathfrak{U}t^n-\mathfrak{T}u^n).
\]
\SourceLeaf{196}
Das lässt sich leicht durch Einsetzen bestätigen. Dass diese Werte, die wir kurz mit \(=\tau,\upsilon\) bezeichnen, stets \textit{ganze Zahlen} sind, zeigen wir wie folgt. Ist \((M,N,P)\) eine Form mit der Determinante \(D\) und \(m\) ein gemeinsamer Teiler von \(M,2N,P\), so sind sowohl \(\mathfrak{T}+N\mathfrak{U}\) als auch \(t^n+Nu^n\) durch \(m\) teilbar und daher auch \(\mathfrak{U}(t^n+Nu^n)-u^n(\mathfrak{T}+N\mathfrak{U})\) oder \(\mathfrak{U}t^n-\mathfrak{T}u^n\). Somit ist \(\upsilon\) ganzzahlig und daher auch \(\tau\), weil \(\tau\tau=D\upsilon\upsilon+mm\).

\(2^\circ\) Offenbar kann \(\upsilon\) nicht \(=0\) sein, denn daraus würde folgen
\[
\mathfrak{UU}t^n t^n=\mathfrak{TT}u^n u^n
\]
oder
\[
\mathfrak{UU}(Du^n u^n+mm)=u^n u^n(D\mathfrak{UU}+mm)
\]
oder \(\mathfrak{UU}=u^n u^n\), im Widerspruch zur Voraussetzung \(\mathfrak{U}>u^n\). Da also außer dem Wert \(0\) der kleinste Wert von \(u\) gleich \(U\) ist, ist \(\upsilon\) sicher nicht kleiner als \(U\).

\(3^\circ\) Aus den Werten von \(t^n,t^{n+1},u^n,u^{n+1}\) lässt sich leicht bestätigen, dass
\[
mU=u^{n+1}t^n-t^{n+1}u^n.
\]
Daher ist \(\mathfrak{U}t^n-\mathfrak{T}u^n\) sicher nicht kleiner als \(u^{n+1}t^n-t^{n+1}u^n\).

\(4^\circ\) Nun ergibt die Gleichung \(\mathfrak{TT}-D\mathfrak{UU}=mm\)
\[
\frac{\mathfrak{T}}{\mathfrak{U}}=\surd\left(D+\frac{mm}{\mathfrak{UU}}\right)
\]
und entsprechend
\[
\frac{t^{n+1}}{u^{n+1}}=\surd\left(D+\frac{mm}{u^{n+1}u^{n+1}}\right)
\]
woraus leicht \(\frac{\mathfrak{T}}{\mathfrak{U}}>\frac{t^{n+1}}{u^{n+1}}\) folgt. Daraus und aus dem Schluss in \(3^\circ\) folgt
\[
\begin{aligned}&(\mathfrak{U}t^n-\mathfrak{T}u^n)\left(t^n+u^n\frac{\mathfrak{T}}{\mathfrak{U}}\right)\\&\qquad>(u^{n+1}t^n-t^{n+1}u^n)\left(t^n+u^n\frac{t^{n+1}}{u^{n+1}}\right)\end{aligned}
\]
oder, wenn wir ausmultiplizieren und für \(\mathfrak{TT},\ t^n t^n,\ t^{n+1}t^{n+1}\) ihre Werte \(D\mathfrak{UU}+mm,\ Du^n u^n+mm,\ Du^{n+1}u^{n+1}+mm\) einsetzen,
\[
\frac1{\mathfrak{U}}(\mathfrak{UU}-u^n u^n)>\frac1{u^{n+1}}(u^{n+1}u^{n+1}-u^n u^n)
\]
Da beide Größen offenbar positiv sind, ergibt sich durch Umstellen \(\mathfrak{U}+\frac{u^n u^n}{u^{n+1}}>u^{n+1}+\frac{u^n u^n}{\mathfrak{U}}\), \textit{ein Widerspruch}, weil der erste Teil der ersten Größe \textit{kleiner} ist
\SourceLeaf{197}
als der erste Teil der zweiten und ebenso ihr zweiter Teil kleiner als der zweite Teil der anderen. Daher kann die Annahme nicht gelten, und die Folgen \(t^0,t',t''\) usw., \(u^0,u',u''\) usw. stellen sämtliche positiven Werte von \(t,u\) dar.

\textit{Beispiel.} Für \(D=61,\ m=2\) fanden wir als kleinste positive Werte von \(t,u\) die Werte \(1523,195\). Daher stellen die folgenden Formeln \textit{alle} positiven Werte dar:
\[
\begin{aligned}
t&=\left(\frac{1523}{2}+\frac{195}{2}\surd61\right)^e+\left(\frac{1523}{2}-\frac{195}{2}\surd61\right)^e\\
u&=\frac1{\surd61}\left(\left(\frac{1523}{2}+\frac{195}{2}\surd61\right)^e-\left(\frac{1523}{2}-\frac{195}{2}\surd61\right)^e\right)
\end{aligned}
\]
Wir finden
\[
\begin{aligned}
t^0&=2,\quad t'=1523,\\
t''&=1523t'-t^0=2319527,\\
t'''&=1523t''-t'=3532618098\text{ usw.}\\
u^0&=0,\quad u'=195,\\
u''&=1523u'-u^0=296985,\\
u'''&=1523u''-u'=452307960\text{ usw.}
\end{aligned}
\]
